%%
%% part2-book-components/ch07-text.tex
%%
%% Chapter 7: Text and Paragraphs — How to write Persian
%% text, mix it with Latin text, and control paragraph
%% formatting.
%%

\chapter{متن و پاراگراف}
\label{chap:text}

\section{چرا این موضوع مهم است؟}
\label{sec:text-why}

متن، بخش اصلی هر کتاب است. در کتاب‌های فنی
فارسی، متن ترکیبی از فارسی و لاتین است و
مدیریت درست این ترکیب، کیفیت نهایی کتاب را
تعیین می‌کند.

\lr{MatinBook} با استفاده از \lr{xepersian} و
تنظیمات دقیق تایپوگرافی، این ترکیب را به‌درستی
مدیریت می‌کند. هدف این فصل، آموزش نحوه‌ی نوشتن
متن فارسی، ترکیب آن با لاتین، و کنترل پاراگراف‌ها
است.

\mnote{تایپوگرافی خوب، خوانایی را افزایش می‌دهد و خستگی چشم خواننده را کاهش می‌دهد.}

\section{متن فارسی}
\label{sec:text-persian}

نوشتن متن فارسی در \lr{MatinBook}، ساده‌ترین
کار ممکن است: فقط متن را به فارسی بنویسید و
\lr{LaTeX} بقیه‌ی کارها را انجام می‌دهد.

\begin{latin}
\begin{minted}{latex}
|\pc{این یک متن فارسی ساده است. لاتک به‌طور خودکار}|%
|\pc{جهت متن را راست‌به‌چپ تشخیص می‌دهد.}|%
\end{minted}
\end{latin}

\mnote{نیازی به هیچ دستور خاصی برای متن فارسی نیست. فقط متن را بنویسید.}

\subsection{فاصله‌گذاری در فارسی}
\label{sec:text-persian-spacing}

در فارسی، فاصله‌ی میان کلمات و نشانه‌های نگارشی،
قواعد خاصی دارد:

\begin{itemize}
    \item \textbf{فاصله‌ی معمولی:} بین کلمات.
    \item \textbf{نیم‌فاصله (\lr{ZWNJ}):} در کلماتی
    مانند «می‌رود»، «کتاب‌ها»، «برنامه‌نویسی».
    \item \textbf{فاصله‌ی نشکن (\lr{\textasciitilde}):}
    برای جلوگیری از شکستن خط.
\end{itemize}

\mnote{نیم‌فاصله را می‌توانید با \lr{Ctrl+Shift+2} در ویندوز یا \lr{Shift+Space} در لینوکس تایپ کنید.}

\subsection{نشانه‌های نگارشی فارسی}
\label{sec:text-persian-punctuation}

\lr{MatinBook} از تمام نشانه‌های نگارشی فارسی
پشتیبانی می‌کند:

\begin{itemize}
    \item \textbf{گیومه:} «نقل قول مستقیم».
    \item \textbf{گیومه‌ی تکی:} ‹نقل قول درونی›.
    \item \textbf{پرانتز:} (متن داخل پرانتز).
    \item \textbf{کروشه:} [متن داخل کروشه].
    \item \textbf{نقطه‌ویرگول:} نشانه‌ی مکث؛ ادامه.
    \item \textbf{ویرگول:} جداسازی، ادامه.
    \item \textbf{علامت سؤال:} چرا؟
    \item \textbf{علامت تعجب:} شگفت‌انگیز!
\end{itemize}

\mnote{در فارسی، نقطه‌ویرگول «؛» و ویرگول «،» با نسخه‌ی لاتین متفاوت هستند. \lr{xepersian} این تفاوت را به‌درستی مدیریت می‌کند.}

\section{ترکیب فارسی و لاتین}
\label{sec:text-latin}

در کتاب‌های فنی، ترکیب متن فارسی با واژگان
لاتین امری اجتناب‌ناپذیر است. \lr{MatinBook}
سه روش برای این ترکیب ارائه می‌دهد:

\begin{enumerate}
    \item \cmd{lr\{...\}}: برای واژگان و
    عبارات کوتاه لاتین در متن فارسی.
    \item \cmd{rl\{...\}}: برای متن فارسی در
    محیط لاتین (مخصوص \lr{tikzpicture}).
    \item \cmd{begin\{latin\}}: برای پاراگراف‌های
    کامل لاتین.
\end{enumerate}

\subsection{واژگان لاتین در متن فارسی}
\label{sec:text-latin-inline}

برای واژگان و عبارات کوتاه لاتین، از دستور
\cmd{lr\{...\}} استفاده کنید:

\begin{latin}
\begin{minted}{latex}
|\pc{این کتاب درباره}|\lr{machine learning}|\pc{و}|\lr{deep learning}|\pc{است.}|%
\end{minted}
\end{latin}

خروجی: این کتاب درباره \lr{machine learning}
و \lr{deep learning} است.

\mnote{فراموش کردن \lr{\textbackslash lr} باعث می‌شود واژه‌ی لاتین با متن فارسی قاطی شود و به‌صورت نامفهوم نمایش داده شود.}

\subsection{متن فارسی در محیط لاتین}
\label{sec:text-persian-in-latin}

برای متن فارسی در محیط لاتین (مخصوص
\lr{tikzpicture})، از دستور \cmd{rl\{...\}}
استفاده کنید:

\begin{latin}
\begin{minted}{latex}
\begin{tikzpicture}
    \node {|\rl{متن فارسی}|};
\end{tikzpicture}
\end{minted}
\end{latin}

\mnote{در \lr{tikzpicture}، جهت متن به‌طور خودکار لاتین است. برای متن فارسی، از \lr{\textbackslash rl} استفاده کنید.}

\subsection{پاراگراف‌های کامل لاتین}
\label{sec:text-latin-block}

برای پاراگراف‌های کامل لاتین (مثلاً یک متن
انگلیسی چند خطی)، از محیط \env{latin} استفاده
کنید:

\begin{latin}
\begin{minted}{latex}
\begin{latin}
This is a complete English paragraph.
It can span multiple lines and will be
typeset left-to-right.
\end{latin}
\end{minted}
\end{latin}

\mnote{محیط \lr{latin} برای کد، الگوریتم، نمودار و هر متن لاتین طولانی مناسب است.}

\subsection{جدول تصمیم‌گیری}
\label{sec:text-decision-table}

جدول \cref{tab:latin-decision} به شما کمک می‌کند
که در هر موقعیت، از کدام دستور استفاده کنید.

\begin{matintable}{انتخاب دستور مناسب برای ترکیب فارسی و لاتین}{tab:latin-decision}
\begin{tblr}{
    width = \textwidth,
    colspec = {X[l] X[c]},
    hlines, vlines,
    colsep = 8pt,
    rowsep = 4pt,
    row{1} = {font=\bfseries},
}
موقعیت & دستور \\
واژه‌ی لاتین در متن فارسی & \lr{\textbackslash lr\{...\}} \\
عبارت کوتاه لاتین & \lr{\textbackslash lr\{...\}} \\
پاراگراف کامل لاتین & \lr{\textbackslash begin\{latin\}} \\
متن فارسی در \lr{tikzpicture} & \lr{\textbackslash rl\{...\}} \\
کد برنامه‌نویسی & \lr{\textbackslash begin\{latin\}} + \lr{minted} \\
الگوریتم & \lr{\textbackslash begin\{latin\}} + \lr{algorithm} \\
فرمول ریاضی & \lr{\textbackslash lr\{...\}} در متن فارسی \\
\lr{URL} & \lr{\textbackslash lr\{...\}} \\
\end{tblr}
\end{matintable}

\mnote{اگر شک دارید، از \lr{\textbackslash lr} استفاده کنید. این دستور در ۹۰٪ موارد کار می‌کند.}

\section{پاراگراف‌ها}
\label{sec:text-paragraphs}

پاراگراف، واحد اصلی متن است. \lr{MatinBook}
تنظیمات زیر را برای پاراگراف‌ها اعمال می‌کند:

\begin{itemize}
    \item \textbf{تورفتگی:} \pnum{۱em} برای
    پاراگراف‌های دوم به بعد.
    \item \textbf{فاصله‌ی بین پاراگراف‌ها:}
    \pnum{۰pt} (بدون فاصله‌ی اضافی).
    \item \textbf{فاصله‌ی خطوط:} \pnum{۱.۱۵}.
\end{itemize}

\mnote{پاراگراف اول هر بخش، تورفتگی ندارد. این یک استاندارد در حروف‌چینی فارسی است.}

\subsection{شکستن پاراگراف}
\label{sec:text-paragraph-break}

برای شکستن پاراگراف، یک خط خالی بین دو متن
قرار دهید:

\begin{latin}
\begin{minted}{latex}
|\pc{این پاراگراف اول است.}|%

|\pc{این پاراگراف دوم است.}|%
\end{minted}
\end{latin}

\mnote{هرگز از \lr{\textbackslash\textbackslash} برای شکستن پاراگراف استفاده نکنید. این دستور برای شکستن خط در یک پاراگراف است، نه برای ایجاد پاراگراف جدید.}

\subsection{طول پاراگراف}
\label{sec:text-paragraph-length}

طبق اصول نگارش حرفه‌ای:

\begin{itemize}
    \item \textbf{طول ایده‌آل:} ۳ تا ۷ جمله.
    \item \textbf{حداکثر:} ۱۰ جمله.
    \item \textbf{حداقل:} ۲ جمله.
\end{itemize}

\mnote{اگر پاراگراف بیش از ۱۰ جمله است، آن را به دو پاراگراف تقسیم کنید. اگر فقط یک جمله است، آن را با پاراگراف بعدی ادغام کنید.}

\section{متن تأکیدی}
\label{sec:text-emphasis}

\lr{MatinBook} از دستورات استاندارد \lr{LaTeX}
برای تأکید پشتیبانی می‌کند. جدول
\cref{tab:emphasis-commands} این دستورات را
خلاصه می‌کند:

\begin{matintable}{دستورات تأکید در متن}{tab:emphasis-commands}
\begin{tblr}{
    width = 0.7\textwidth,
    colspec = {l X[c]},
    hlines, vlines,
    colsep = 8pt,
    rowsep = 4pt,
    row{1} = {font=\bfseries},
}
نوع تأکید & دستور \\
پررنگ & \lr{\textbackslash textbf\{...\}} \\
مورب & \lr{\textbackslash textit\{...\}} \\
زیرخط‌دار & \lr{\textbackslash underline\{...\}} \\
مونواسپیس & \lr{\textbackslash texttt\{...\}} \\
\end{tblr}
\end{matintable}

در ادامه، هر یک را با یک مثال نشان می‌دهیم.

\subsection{متن پررنگ}

برای پررنگ کردن متن، از دستور \cmd{textbf}
استفاده کنید. مثال:

\begin{quote}
    این متن \textbf{پررنگ} است.
\end{quote}

\mnote{در متن فارسی، \lr{\textbackslash textbf} به‌درستی کار می‌کند و متن را پررنگ نمایش می‌دهد.}

\subsection{متن مورب}

برای مورب کردن متن، از دستور \cmd{textit}
استفاده کنید. مثال:

\begin{quote}
    این متن \textit{مورب} است.
\end{quote}

\mnote{متن مورب در فارسی کمتر استفاده می‌شود، چون خوانایی را کاهش می‌دهد. برای تأکید، از پررنگ استفاده کنید.}

\subsection{متن زیرخط‌دار}

برای زیرخط‌دار کردن متن، از دستور \cmd{underline}
استفاده کنید. مثال:

\begin{quote}
    این متن \underline{زیرخط‌دار} است.
\end{quote}

\mnote{استفاده از زیرخط در متن رسمی توصیه نمی‌شود، چون با لینک‌های \lr{hyperref} اشتباه می‌شود.}

\subsection{متن مونواسپیس}

برای نمایش متن با فونت مونواسپیس، از دستور
\cmd{texttt} استفاده کنید. مثال:

\begin{quote}
    این متن \texttt{monospace} است.
\end{quote}

\mnote{برای متن مونواسپیس فارسی، از \lr{\textbackslash texttt} استفاده نکنید، چون فونت مونواسپیس ممکن است فارسی را پشتیبانی نکند.}

\section{فهرست‌ها}
\label{sec:text-lists}

\lr{MatinBook} از سه نوع فهرست پشتیبانی می‌کند.
هر یک را با یک مثال نشان می‌دهیم.

\subsection{فهرست نشانه‌دار}

فهرست نشانه‌دار، پرکاربردترین نوع فهرست است.
مثال:

\begin{quote}
    \begin{itemize}
        \item مورد اول.
        \item مورد دوم.
        \item مورد سوم.
    \end{itemize}
\end{quote}

\mnote{از فهرست نشانه‌دار برای موارد بدون ترتیب مشخص استفاده کنید.}

\subsection{فهرست شماره‌دار}

فهرست شماره‌دار، برای مواردی که ترتیب مهم
است، استفاده می‌شود. مثال:

\begin{quote}
    \begin{enumerate}
        \item مورد اول.
        \item مورد دوم.
        \item مورد سوم.
    \end{enumerate}
\end{quote}

\mnote{از فهرست شماره‌دار برای مراحل، دستورالعمل‌ها یا اولویت‌ها استفاده کنید.}

\subsection{فهرست توصیفی}

فهرست توصیفی، برای تعریف اصطلاحات استفاده
می‌شود. مثال:

\begin{quote}
    \begin{description}
        \item[اصطلاح اول:] توضیح اول.
        \item[اصطلاح دوم:] توضیح دوم.
    \end{description}
\end{quote}

\mnote{از فهرست توصیفی برای واژه‌نامه یا تعریف اصطلاحات استفاده کنید.}

\mnote{در همه‌ی فهرست‌ها، آیتم‌ها باید ساختار گرامری یکسان داشته باشند. این اصل «ساختار موازی» نام دارد.}

\section{خطاهای رایج}
\label{sec:text-mistakes}

\subsection{فراموش کردن \cmd{lr}}
\label{sec:text-mistake-lr}

اگر \cmd{lr} را برای واژگان لاتین فراموش کنید،
نتیجه به‌صورت نامفهوم نمایش داده می‌شود:

\begin{latin}
\begin{minted}{latex}
% |\pc{اشتباه}|:
|\pc{این کتاب درباره}|machine learning|\pc{است.}|%

% |\pc{درست}|:
|\pc{این کتاب درباره}|\lr{machine learning}|\pc{است.}|%
\end{minted}
\end{latin}

\subsection{استفاده از \cmd{\textbackslash\textbackslash} برای پاراگراف}
\label{sec:text-mistake-parbreak}

دستور \cmd{\textbackslash\textbackslash} برای شکستن
خط در یک پاراگراف است، نه برای ایجاد پاراگراف
جدید. برای پاراگراف جدید، از یک خط خالی
استفاده کنید.

\subsection{پاراگراف بسیار طولانی}
\label{sec:text-mistake-longpara}

پاراگراف‌های بیش از ۱۰ جمله، خواننده را خسته
می‌کنند. آن‌ها را به پاراگراف‌های کوچک‌تر
تقسیم کنید.

\section{تمرین‌ها}
\label{sec:text-exercises}

\begin{exercise}
    یک پاراگراف فارسی با پنج جمله بنویسید
    که شامل سه واژه‌ی لاتین باشد. از
    \cmd{lr} برای واژگان لاتین استفاده کنید.
    \label{exc:text-paragraph}
\end{exercise}

\begin{exercise}
    یک فهرست نشانه‌دار با پنج مورد بنویسید
    که همه‌ی موارد ساختار گرامری یکسان داشته
    باشند.
    \label{exc:text-list}
\end{exercise}

\begin{exercise}
    یک پاراگراف لاتین با محیط \env{latin}
    بنویسید و آن را با متن فارسی ترکیب کنید.
    \label{exc:text-latin}
\end{exercise}

\begin{exercise}
    یک متن فارسی با نشانه‌های نگارشی مختلف
    (گیومه، پرانتز، کروشه، نقطه‌ویرگول) بنویسید.
    \label{exc:text-punctuation}
\end{exercise}

\section{خلاصه}
\label{sec:text-summary}

در این فصل، با متن و پاراگراف در \lr{MatinBook}
آشنا شدیم:

\begin{itemize}
    \item متن فارسی به‌طور ساده نوشته می‌شود و
    \lr{LaTeX} جهت را خودکار تشخیص می‌دهد.

    \item برای واژگان لاتین در متن فارسی، از
    \cmd{lr\{...\}} استفاده کنید.

    \item برای متن فارسی در محیط لاتین، از
    \cmd{rl\{...\}} استفاده کنید.

    \item برای پاراگراف‌های کامل لاتین، از
    محیط \env{latin} استفاده کنید.

    \item پاراگراف‌ها با یک خط خالی از هم جدا
    می‌شوند.

    \item طول ایده‌آل پاراگراف، ۳ تا ۷ جمله
    است.

    \item خطاهای رایج شامل فراموش کردن
    \cmd{lr}، استفاده از \cmd{\textbackslash\textbackslash}
    برای پاراگراف، و پاراگراف‌های بسیار طولانی
    است.
\end{itemize}

در فصل بعدی (\cref{chap:environments})، با
محیط‌های علمی (قضیه، لم، تعریف، مثال، اثبات)
آشنا می‌شوید و یاد می‌گیرید که چگونه از آن‌ها
استفاده کنید.